\begin{table}[!htbp]
\centering
\caption{\emph{Checklist} pemenuhan kebutuhan fungsional F1--F6}
\label{tab:checklist-f1-f6}
\small
\renewcommand{\arraystretch}{1.2}

\begin{tabular}{|p{0.9cm}|>{\raggedright\arraybackslash}p{5.0cm}|>{\raggedright\arraybackslash}p{6.9cm}|}
\hline
\textbf{Req.} &
\textbf{Komponen Implementasi} &
\textbf{Kriteria Verifikasi} \\
\hline

F1 &
\emph{Pipeline} konstruksi dan densifikasi KG berbobot dari SORD &
KG terbangun dari seluruh populasi SORD, bobot simpul dan sisi terisi dari sinyal komunitas, dan proporsi simpul terisolasi $\leq 30\%$ setelah densifikasi \\
\hline

F2 &
Indeks vektor FAISS untuk pertanyaan dan untuk jawaban (\texttt{all-MiniLM-L6-v2}) &
Kedua indeks dapat di-\emph{query} dan mengembalikan pertanyaan atau jawaban paling mirip beserta skor kemiripannya \\
\hline

F3 &
\emph{Entity anchoring} pada tingkat jawaban berbasis pencarian vektor &
Setiap kueri memperoleh \emph{seed nodes} di KG dari pertanyaan induk jawaban yang paling mirip \\
\hline

F4 &
\emph{Hybrid retrieval} (\emph{anchoring}, \emph{traversal}, dan \emph{semantic expansion}) serta fusi berbobot pada kedua varian Kondisi C &
Kandidat dari ketiga sumber digabung dalam satu \emph{pool}, diberi bobot jalur menurut definisi bobot jawaban masing-masing varian, disaring (jawaban tidak diterima dan bernilai negatif dikeluarkan sebelum pemilihan), lalu $k$ kandidat teratas disusun menjadi \emph{unified context window} yang deterministik pada eksekusi berulang \\
\hline

F5 &
\emph{Grounded generation} dengan \emph{dual-constraint} &
Output menyertakan jawaban, kutipan \texttt{[SO-<id>]} yang tervalidasi terhadap konteks, dan catatan \emph{provenance} per item konteks \\
\hline

F6 &
\emph{Pipeline} evaluasi empat kondisi &
Kondisi A, B, C-Acc, dan C-Acc+Vote dapat dijalankan pada sampel yang sama dengan konfigurasi, \emph{prompt}, dan pengecualian kebocoran yang tercatat \\
\hline

\end{tabular}
\end{table}